% Quarantined v4 research proof.  The finite-residue-field step is
% deliberate: no equality of analytic interiors is asserted after an
% arbitrary complete scalar extension.

\begin{proof}[Proof of
\cref{prop:p062-strict-product-fibre-slicing-v4-research}]
Write \(\pi\colon X\to B\) for the projection.  It is flat,
boundaryless, and of pure relative dimension one.  Fix \(x\in B\) and a
compactly supported ordinary CLD-smooth function \(\varphi\) on
\(
 X_x=Z_{\mathscr H(x)}
\).

Since \(B\) is good and strict, choose a strict affinoid neighbourhood
\(U_0\) of \(x\) in \(B\) such that
\(
 x\in\operatorname{Int}(U_0/B)
\).  Since \(B\) is locally compact and Hausdorff, there is an ordinary
open neighbourhood \(V\) of \(x\) whose closure is compact and contained
in \(\operatorname{Int}(U_0/B)\).  Put
\[
                         Y_0=U_0\times_KZ,
                         \qquad F=(Y_0)_x=X_x.
\]
The space \(Y_0\) is good and topologically separated, hence almost good
and locally holomorphically separated.  Moreover,
\begin{equation}\label{eq:p062-endpoint-relative-interior-support-v4-research}
 V\times_KZ\subseteq\operatorname{Int}(Y_0/X),
\end{equation}
because \(V\times_KZ\) is an ordinary open subset of \(X\) contained in
the analytic domain \(Y_0\); here we use the identification of relative
interior with topological interior for an analytic domain
\cite[Proposition~1.5.5(i)]{Berk93}.  We spread \(\varphi\) to a smooth
function on \(Y_0\), exactly on the fibre, and keep its support inside this
ordinary open subset of \(X\).

Fix \(y\in\operatorname{supp}(\varphi)\).  The ordinary-open local moment
presentation of an ordinary form
\cite[Proposition~15.1.1(b)]{CLD25} gives a neighbourhood \(V_y'\) of \(y\) in
\(F\), invertible \(\mathscr H(x)\)-analytic functions
\(f_{y,1},\ldots,f_{y,N_y}\) there, and a smooth real function \(\rho_y\)
such that
\[
 \varphi=\rho_y\bigl(-\log|f_{y,1}|,\ldots,
                         -\log|f_{y,N_y}|\bigr)
 \qquad\text{on }V_y'.
\]
The completed residue field of \(y\), computed in the analytic fibre
\(F\), is canonically the same as that computed in \(Y_0\).  Indeed, on
affinoid charts associated with \(A\to C\), adjoining the image of
\(\mathscr H(x)\) to \(C\) does not change the completion of the fraction
field of the image of \(C\): the latter is already dense in
\(\mathscr H_{Y_0}(y)\), and the image of \(\mathscr H(x)\) lies in that
closure.

Since \(Y_0\) is good at \(y\), its local residue field is dense in
\(\mathscr H_{Y_0}(y)\) by \cite[\S11.1.10]{CLD25}.  For each \(k\), choose
a local ambient meromorphic residue class represented by an analytic germ
\(g_{y,k}\in\mathcal O_{Y_0,y}\) such that
\begin{equation}\label{eq:p062-endpoint-moment-approximation-v4-research}
             |g_{y,k}(y)-f_{y,k}(y)|<|f_{y,k}(y)|.
\end{equation}
After shrinking around \(y\), every \(g_{y,k}\) is an analytic unit and
the ultrametric inequality gives the literal identities
\begin{equation}\label{eq:p062-endpoint-moment-equality-v4-research}
                    |g_{y,k}|_F=|f_{y,k}|.
\end{equation}
We also shrink the ambient neighbourhood so that it is contained in
\(V\times_KZ\).  Thus the ordinary smooth ambient function
\[
 \Phi_y=\rho_y\bigl(-\log|g_{y,1}|,\ldots,
                         -\log|g_{y,N_y}|\bigr)
\]
restricts exactly to \(\varphi\) on a fibre neighbourhood of \(y\), and is
defined on an ambient neighbourhood contained in \(V\times_KZ\).

The support of \(\varphi\) is compact.  Choose finitely many such ambient
neighbourhoods \(W_1,\ldots,W_N\), all contained in \(V\times_KZ\), and use
\cite[Corollary~15.2.9]{CLD25} to choose nonnegative ordinary smooth
functions \(\theta_i\) on \(Y_0\), compactly supported in \(W_i\), whose
sum is one near \(\operatorname{supp}(\varphi)\).  Each
\(\theta_i\Phi_i\) extends smoothly by zero on \(Y_0\).  The sum
\[
                              \Phi=\sum_i\theta_i\Phi_i
\]
has compact support contained in \(V\times_KZ\) and satisfies
\begin{equation}\label{eq:p062-endpoint-exact-spread-v4-research}
                              \Phi|_F=\varphi.
\end{equation}
By \eqref{eq:p062-endpoint-relative-interior-support-v4-research}, this
compact support is contained in the ordinary interior of \(Y_0\) relative
to \(X\).  We may therefore regard \(\Phi\) as a compactly supported smooth
function on the ordinary open subspace \(V\times_KZ\); extension by zero
to \(X\), as in \cite[Proposition~6.4.5(a)]{CLD25}, is ordinary CLD-smooth.
We use the same symbol for this extension.

Set
\[
                         \tau=\mathrm d'\mathrm d''\Phi.
\]
The total space \(X\) is good and topologically separated, hence locally
holomorphically separated.  Decompose the compactly supported symmetric
form \(\tau\) as
\(
 \tau=\tau^+-\tau^-
\)
with \(\tau^\pm\) positive and compactly supported, using
\cite[Proposition~6.5.2(b)]{CLD25}.  Let
\(
 C=\operatorname{supp}(\tau^+)\cup\operatorname{supp}(\tau^-)
\), a compact subset of \(X\).  By
\cite[Corollary~15.2.5]{CLD25}, choose an ordinary smooth compactly
supported function \(\chi\) on \(X\) that equals one on a neighbourhood of
\(C\).  The function \(\chi h\) is continuous and compactly supported.
Applying \cite[Theorem~15.2.6]{CLD25} with errors tending to zero gives
ordinary smooth compactly supported functions \(h_\ell\) such that
\begin{equation}\label{eq:p062-endpoint-cutoff-approximation-v4-research}
 \lVert h_\ell-\chi h\rVert_{C^0(X)}\longrightarrow0.
\end{equation}
In particular, \(h_\ell\to h\) uniformly on \(C\).

For \(t\in B\), set
\[
                 G(t)=\int_{X_t}h_t\,\tau|_{X_t}.
\]
Fibre integration makes
\(
 t\mapsto\int_{X_t}\tau^\pm
\)
continuous by \cite[Theorem~16.3.3(d)]{CLD25}.  These functions vanish
away from the compact projections of the supports of \(\tau^\pm\), and are
therefore bounded on \(B\).  Consequently
\begin{equation}\label{eq:p062-endpoint-uniform-fibre-limit-v4-research}
 \sup_{t\in B}\left|
   \int_{X_t}(h_\ell-h)\tau\right|
 \leq \lVert h_\ell-h\rVert_{C^0(C)}
       \sup_{t\in B}\left(\int_{X_t}\tau^++\int_{X_t}\tau^-\right)
 \longrightarrow0.
\end{equation}
The smooth fibre integrals are continuous, so \(G\) is continuous.  The
same limit, combined with the projection formula
\cite[Theorem~16.5.4 and Corollary~16.5.5]{CLD25}, gives, for every
compactly supported smooth top-degree form \(\alpha\) on \(B\),
\begin{equation}\label{eq:p062-endpoint-continuous-fubini-v4-research}
 \int_X h\,\tau\wedge\pi^*\alpha
                         =\int_B G\,\alpha.
\end{equation}
Indeed, the projection formula applies to each smooth form
\(h_\ell\tau\); the fibre side converges uniformly by
\eqref{eq:p062-endpoint-uniform-fibre-limit-v4-research}, and the total
space side converges against the fixed compactly supported signed
top-degree measure \(\tau\wedge\pi^*\alpha\).

Since \(\alpha\) has top bidegree on the \(m\)-dimensional base,
\(
 \mathrm d'\alpha=\mathrm d''\alpha=0
\), including when \(m=0\).  Therefore
\[
 \mathrm d'\mathrm d''(\Phi\pi^*\alpha)
       =(\mathrm d'\mathrm d''\Phi)\wedge\pi^*\alpha.
\]
The function \(\Phi\) is now a global compactly supported smooth function,
so \(\Phi\pi^*\alpha\) is a legitimate compactly supported test form on
\(X\).  Pairing
\eqref{eq:p062-general-product-zero-current-v4-research} with this test form
and using \eqref{eq:p062-endpoint-continuous-fubini-v4-research} gives
\(
 \int_B G\alpha=0
\)
for every \(\alpha\).  The strict pure-dimensional space \(B\) is very
rich by \cite[Theorem~16.4.4(b)]{CLD25}; hence continuous functions inject
into \((0,0)\)-currents by \cite[Corollary~7.3.12]{CLD25}.  Thus \(G=0\).
At \(x\), \eqref{eq:p062-endpoint-exact-spread-v4-research} now yields
\[
 0=G(x)=\int_{X_x}h_x\,\mathrm d'\mathrm d''\varphi
       =\bigl\langle\mathrm d'\mathrm d''[h_x],\varphi\bigr\rangle.
\]
This proves \eqref{eq:p062-general-fibre-zero-current-v4-research}.
\end{proof}

\begin{proof}[Proof of \cref{thm:p062-endpoint-polycylinder-v4-research}]
If \(p=q=0\), then \(P=\mathcal M(K)\), and the assertion holds with
\(c=h(*)\).  We henceforth assume \(p+q>0\).

We separate the absolute-interior issue from the affine interpolation.  For
every factor \(F_\nu\) of \(P\), write
\(
 F_\nu^\natural=\operatorname{Int}(F_\nu)
\)
for its absolute analytic interior.  Thus an open factor is unchanged; a
closed disc loses only its endpoint Gauss point; and an annulus loses the
Gauss point at each included endpoint but retains every nonskeletal
endpoint residue branch.  The latter one-dimensional description is the
topological input in the verified endpoint theorem
\cref{prop:p062-annulus-endpoint-arbitrary-field}.

Let \(M\) be the product of the ambient affine lines and analytic tori, and
put
\begin{equation}\label{eq:p062-endpoint-coordinate-interior-v4-research}
 Q=\bigcap_\nu\operatorname{pr}_\nu^{-1}(F_\nu^\natural)\subset P.
\end{equation}
Every \(F_\nu^\natural\) is an ordinary open subset of its ambient curve.
Hence \(Q\) is an ordinary open subset of \(M\) contained in \(P\).  It is
therefore contained in the topological interior of the analytic domain
\(P\subset M\).  Berkovich's relative-interior theorem and composition
formula \cite[Proposition~1.5.5(i)--(ii)]{Berk93}, together with the
boundarylessness of \(M\), give
\begin{equation}\label{eq:p062-endpoint-Q-in-interior-v4-research}
                              Q\subset\operatorname{Int}(P).
\end{equation}
Every rigid point of \(P\) lies in \(Q\): its coordinate images are rigid,
whereas the excluded one-dimensional endpoint Gauss points are nonrigid.
Rigid points are dense locally in a strict analytic space
\cite[Proposition~2.1.15]{Berk12}.  Consequently
\begin{equation}\label{eq:p062-endpoint-Q-dense-v4-research}
                              \overline Q=P.
\end{equation}

We next record why slicing at a rigid base point sees the full
one-dimensional absolute interior after passage to the residue field.  Fix
a factor \(F_\nu\), let
\[
 B_\nu=\prod_{\mu\ne\nu}F_\mu^\natural,
 \qquad Z_\nu=F_\nu^\natural.
\]
Here \(B_\nu=\mathcal M(K)\) when the displayed product is empty.  The
product \(B_\nu\times_KZ_\nu\) is \(Q\), so the current equation holds
there by \eqref{eq:p062-endpoint-Q-in-interior-v4-research}.  The spaces
\(B_\nu\) and \(Z_\nu\) are good, topologically separated, pure and strict,
and \(Z_\nu\) is boundaryless.  Thus
\cref{prop:p062-strict-product-fibre-slicing-v4-research} applies, including
when \(B_\nu\) has dimension zero.

Let \(x\in B_\nu\) be rigid and set \(E=\mathscr H(x)\).  Then \(E/K\) is
finite.  For the scalar-extension map
\(
 p_E\colon F_{\nu,E}\to F_\nu
\), one has
\begin{equation}\label{eq:p062-endpoint-finite-interior-equality-v4-research}
       (F_\nu^\natural)_E
       =\operatorname{Int}(F_{\nu,E}/E).
\end{equation}
Indeed, \(p_E\) is finite and hence boundaryless.  Apply the equality case
of Berkovich's composition formula first to
\[
 F_{\nu,E}\longrightarrow F_\nu\longrightarrow\mathcal M(K)
\]
and then to
\[
 F_{\nu,E}\longrightarrow\mathcal M(E)\longrightarrow\mathcal M(K).
\]
The two finite arrows have empty relative boundary, and the remaining maps
are good and locally separated.  The first composition gives
\(
 \operatorname{Int}(F_{\nu,E}/K)
   =p_E^{-1}(\operatorname{Int}(F_\nu/K))
\), while the second identifies its left side with
\(
 \operatorname{Int}(F_{\nu,E}/E)
\).  This is
\eqref{eq:p062-endpoint-finite-interior-equality-v4-research}.  Notice that
this argument uses the finiteness of \(E/K\); no analogous equality is
claimed for an arbitrary complete scalar extension.

Let \(h_x\) be the restriction of \(h\) to the full fibre
\(F_{\nu,E}\).  The slicing proposition and
\eqref{eq:p062-endpoint-finite-interior-equality-v4-research} show that its
Hessian current vanishes on the absolute analytic interior of that fibre.
The verified arbitrary-field one-dimensional results therefore give
\begin{equation}\label{eq:p062-endpoint-one-fibre-classification-v4-research}
 \begin{cases}
 h_x\text{ is constant},&F_\nu\text{ is a disc},\\
 h_x=a_x\log|S_\nu|+b_x,&F_\nu\text{ is an annulus},
 \end{cases}
\end{equation}
where the disc case is
\cref{prop:p062-arbitrary-field-disc}.  For the annulus case, the fully open
domain \(r_\nu<|S_\nu|<R'_\nu\) is covered by the source-correct, quarantined
strict-open result
\cref{prop:p062-arbitrary-field-annulus-v4-source-correct}; the open--closed,
closed--open, and closed--closed domains are covered by
\cref{prop:p062-annulus-endpoint-arbitrary-field}.  Each proposition is
applied over \(E\), and no endpoint proposition is invoked for the fully open
domain.

For each factor let \(F_\nu^-\) denote its coordinatewise strict core:
\[
\begin{aligned}
 F_\nu^-&=\{|T_\nu|<R_\nu\}
              &&\text{for a disc},\\
 F_\nu^-&=\{r_\nu<|S_\nu|<R'_\nu\}
              &&\text{for an annulus},
\end{aligned}
\]
and put \(P^-=\prod_\nu F_\nu^-\).  The source-correct, quarantined strict-open
polycylinder theorem
\cref{thm:p062-arbitrary-field-polycylinder-v4-source-correct} gives constants
\(c,\lambda_1,\ldots,\lambda_q\) for which
\begin{equation}\label{eq:p062-endpoint-core-formula-v4-research}
             h=L:=c+\sum_{j=1}^q\lambda_j\log|S_j|
             \qquad\text{on }P^-.
\end{equation}

We prove \(h=L\) at every rigid point \(y\in P\) by induction on
\[
 e(y)=\#\{\nu:\operatorname{pr}_\nu(y)\notin F_\nu^-\}.
\]
The coordinate image of a rigid point is rigid.  Hence an index counted by
\(e(y)\) is an included endpoint residue branch, never an endpoint Gauss
point.  The case \(e(y)=0\) is
\eqref{eq:p062-endpoint-core-formula-v4-research}.

Suppose \(e(y)>0\), choose a counted coordinate \(\nu\), and let \(x\) be
the projection of \(y\) to the other factors.  Then \(x\in B_\nu\) is
rigid, so \eqref{eq:p062-endpoint-one-fibre-classification-v4-research}
applies over \(E=\mathscr H(x)\).

If \(F_\nu\) is a disc, replace the \(\nu\)-th coordinate by the zero
section of the fibre.  The resulting point \(y_0\) is rigid,
\(e(y_0)=e(y)-1\), and the first line of
\eqref{eq:p062-endpoint-one-fibre-classification-v4-research} gives
\[
                              h(y)=h(y_0).
\]
The function \(L\) has no disc-coordinate term, so the induction hypothesis
gives
\(
 h(y)=h(y_0)=L(y_0)=L(y)
\).

Suppose instead that \(F_\nu\) is an annulus.  Since the valuation on \(E\)
is nontrivial, the divisible hull of \(|E^\times|\) is dense in
\(\mathbf R_{>0}\).  Choose distinct
\[
       \rho_1,\rho_2\in(r_\nu,R'_\nu)
          \cap |E^{\mathrm{alg},\times}|.
\]
For each \(k\), choose a finite extension \(E_k/E\) and an element of
absolute value \(\rho_k\).  It defines a rigid point \(z_k\) of the annulus
fibre, hence a rigid point \(y_k\in P\) above \(x\).  Both points have
\(e(y_k)=e(y)-1\).  By induction,
\begin{equation}\label{eq:p062-endpoint-two-fibre-values-v4-research}
                         h(y_k)=L(y_k)\qquad(k=1,2).
\end{equation}
On this fibre, both \(h_x\) and \(L|_{F_{\nu,E}}\) are affine functions of
\(\log|S_\nu|\).  The two radii are distinct, so
\eqref{eq:p062-endpoint-two-fibre-values-v4-research} makes those affine
functions identical.  In particular \(h(y)=L(y)\).  This completes the
induction.

Rigid points are dense in \(P\), and both \(h\) and \(L\) are continuous.
Thus \eqref{eq:p062-endpoint-polycylinder-affine-v4-research} holds on all
of \(P\).  If \(q=0\), uniqueness of \(c\) follows by evaluating at any
point of the nonempty space \(P\).  If \(q>0\), restrict to the zero
sections in the disc coordinates and to the standard skeleton of the
strict open polyannulus core.  The functions
\(1,\log|S_1|,\ldots,\log|S_q|\) are affinely independent on its open
logarithmic box, which proves uniqueness of all coefficients.

Finally, every \(\log|S_j|\) is the logarithm of the norm of an analytic
unit.  The displayed formula is therefore ordinary CLD-smooth with zero
smooth Hessian and classically pluriharmonic; the already verified
comparison results give weak smoothness and Gubler--Jell--Rabinoff
harmonicity.
\end{proof}
